\chapter{Cronograma y presupuesto}

El cronograma original se utilizó como punto de partida, pero el seguimiento
se actualizó semanalmente según evidencia producida y bloqueos reales. Al cierre
de esta versión estaban desarrolladas las líneas base de simulación, la
integración de monitoreo, la instrumentación y calibración del robot, la
identificación del actuador, el entrenamiento de las políticas PPO y su
evaluación con el robot suspendido. Permanecían abiertas la seguridad
eléctrica, la migración del bucle de control a la Raspberry~Pi y la evaluación
de la locomoción apoyada.

\section{Cronograma planificado}

La figura~\ref{fig:cronograma-anteproyecto} reproduce el cronograma de 32
semanas del anteproyecto. Frente a él, el modelado y la simulación se
extendieron más de lo previsto, la arquitectura electrónica se resolvió de forma
provisional con una interfaz de E/S para no detener los ensayos, y la
evaluación de estabilidad con apoyo quedó fuera del periodo cubierto por esta
versión.

\begin{figure}[H]
\centering
\resizebox{\textwidth}{!}{%
\begin{ganttchart}[
vgrid,
hgrid,
bar/.style={fill=blue!40},
bar height=0.6
]{1}{32}
\gantttitle{Semanas de trabajo}{32} \\
\gantttitlelist{1,...,32}{1} \\
\ganttbar{Definición del problema}{1}{2} \\
\ganttbar{Formulación de objetivos y alcance}{2}{3} \\
\ganttbar{Estado del arte y revisión bibliográfica}{3}{5} \\
\ganttbar{Redacción del marco teórico y alcance}{4}{5} \\
\ganttbar{Definición de metodología y presupuesto}{5}{6} \\
\ganttbar{Modelado cinemático del robot}{7}{15} \\
\ganttbar{Validación de los modelos}{16}{17} \\
\ganttbar{Diseño de arquitectura electrónica}{16}{18} \\
\ganttbar{Implementación electrónica}{17}{20} \\
\ganttbar{Programación del sistema embebido}{19}{23} \\
\ganttbar{Pruebas individuales de articulaciones}{21}{25} \\
\ganttbar{Desarrollo de algoritmos de locomoción}{23}{29} \\
\ganttbar{Evaluación de estabilidad}{26}{31} \\
\ganttbar{Validación experimental}{27}{31} \\
\ganttbar{Redacción del documento final}{28}{31} \\
\ganttbar{Revisión del asesor}{28}{31} \\
\ganttbar{Entrega del trabajo de grado}{31}{32} \\
\ganttbar{Sustentación}{31}{32}
\end{ganttchart}%
}
\caption{Cronograma del anteproyecto. Fuente: elaboración propia.}
\label{fig:cronograma-anteproyecto}
\end{figure}

\section{Presupuesto del anteproyecto}

El presupuesto del anteproyecto consideró los costos directos de componentes
electrónicos y servomotores, con precios de referencia de tiendas de
electrónica y comercio electrónico en Colombia: valores promedio de kits para
los componentes pequeños y precios de mercado para los elementos más costosos.
Los equipos del laboratorio de robótica (osciloscopio, fuente de laboratorio,
estación de soldadura, computadores y la propia plataforma cuadrúpeda) son
infraestructura institucional, por lo que no se incluyó su costo ni se asignó
un prorrateo; si se adoptara, se calcularía según el número de estudiantes o
proyectos que usan los equipos durante su vida útil.

\begin{table}[H]
\centering
\caption{Costos directos estimados en el anteproyecto.}
\label{tab:costos_materiales}
\begin{tabular}{clcr}
\toprule
\textbf{Ítem} & \textbf{Descripción} & \textbf{Cantidad} & \textbf{Total (COP)} \\
\midrule
1 & Raspberry~Pi~5 de 16~GB & 1 & 1.035.300 \\
2 & Módulo IMU BNO055 de 9 GDL & 1 & 65.000 \\
3 & Sensor de temperatura digital DS18B20 & 12 & 5.500 \\
4 & Servomotores & 12 & 384.000 \\
5 & Conector hembra de 2,54~mm 1$\times$14 & 1 & 2.000 \\
6 & Impresión de PCB & 1 & 50.000 \\
7 & Controlador PWM PCA9685 & 1 & 35.000 \\
8 & Regulador DC--DC LM2596ADJ & 1 & 15.000 \\
\midrule
\multicolumn{3}{r}{\textbf{Suma de los ítems}} & \textbf{1.591.800} \\
\bottomrule
\end{tabular}
\end{table}

El resumen del anteproyecto reportó un total general de 2.029.270~COP, sin
costos prorrateados, que no coincide con la suma de los ítems listados
(1.591.800~COP); la diferencia de 437.470~COP no está desglosada. El valor de
los doce sensores DS18B20 (5.500~COP en total) también requiere aclarar si es
unitario o total.

\section{Conciliación con los componentes utilizados}

La plataforma ejecutada se apartó de esa lista. Se empleó una Raspberry~Pi~4 en
lugar de la Raspberry~Pi~5 presupuestada; los servomotores son MG996R; la PCB y los sensores de temperatura no forman parte
de la plataforma ensayada, y se añadieron
potenciómetros y encoders magnéticos AS5600 con su multiplexor TCA9548A, el
multiplexor analógico CD74HC4067 y la interfaz de E/S. La lista completa, con el
estado de cada componente (en uso, en inventario o por adquirir), está en el
anexo (tabla~\ref{tab:lista-materiales}). El presupuesto anterior se conserva
como antecedente y no se presenta como costo final ejecutado; la conciliación
con los precios de compra reales permanece pendiente.
